\documentclass[11pt,reqno]{amsart}
\usepackage{amsmath,amssymb,amsthm,geometry,mathtools}
\geometry{a4paper,margin=1in}

\newtheorem{theorem}{Theorem}[section]
\newtheorem{lemma}[theorem]{Lemma}
\newtheorem{proposition}[theorem]{Proposition}
\newtheorem{conjecture}[theorem]{Conjecture}

\theoremstyle{remark}
\newtheorem{remark}[theorem]{Remark}

\newcommand{\e}{\mathrm{e}}
\newcommand{\m}{\mathfrak{m}}
\newcommand{\M}{\mathfrak{M}}
\newcommand{\R}{\mathbb{R}}
\newcommand{\Z}{\mathbb{Z}}
\newcommand{\N}{\mathbb{N}}

\title{On the Minor Arc Analysis of Sun's 2-4-6-8 Binomial Sum Conjecture}
\author{Raj123-0}
\date{\today}

\begin{document}

\maketitle

\begin{abstract}
We analyze the representation function $r(n) = \#\{(w,x,y,z) \in \mathbb{N}^4 : n = \binom{w+2}{2} + \binom{x+3}{4} + \binom{y+5}{6} + \binom{z+7}{8}\}$ via the Hardy--Littlewood circle method. While the major arc singular series $\mathfrak{S}(n)$ and singular integral $J(n) \asymp n^{1/24}$ are strictly positive, we demonstrate that standard decoupling and Weyl differencing techniques on the minor arcs exhibit an inherent exponent defect of $\Delta = 7/16$. We formalize the precise bilinear cancellation hypothesis required to complete the proof.
\end{abstract}

\section{Circle Method Formulation}

Let $N \ge 1$ be a large integer, and define the scale parameters $P_k = N^{1/k}$ for $k \in \{2, 4, 6, 8\}$. Define the generating exponential sums:
\begin{equation}
    f_k(\alpha) = \sum_{1 \le u \le P_k} \e\left(\alpha \binom{u+k-1}{k}\right) \quad (k \in \{2, 4, 6, 8\}).
\end{equation}
The representation number $r(n)$ for $n \le N$ is given by the Fourier integral:
\begin{equation}
    r(n) = \int_0^1 f_2(\alpha) f_4(\alpha) f_6(\alpha) f_8(\alpha) \e(-n\alpha) \, d\alpha.
\end{equation}

\section{Major and Minor Arc Dissection}

Let $0 < \theta < 1/8$. For coprime integers $1 \le a \le q \le N^\theta$, define the major arcs:
\begin{equation}
    \M(q, a) = \left\{ \alpha \in [0, 1) : \left| \alpha - \frac{a}{q} \right| \le \frac{1}{q N^{1 - \theta}} \right\}, \qquad \M = \bigcup_{q \le N^\theta} \bigcup_{\substack{a=1 \\ (a,q)=1}}^q \M(q, a).
\end{equation}
The minor arcs are defined as the complement $\m = [0, 1) \setminus \M$.

\begin{proposition}[Major Arc Asymptotics]
For $n \asymp N$, the major arc integral satisfies:
\begin{equation}
    \int_{\M} f_2(\alpha) f_4(\alpha) f_6(\alpha) f_8(\alpha) \e(-n\alpha) \, d\alpha = \mathfrak{S}(n) J(n) + \mathcal{O}\left(N^{1/24 - \varepsilon}\right),
\end{equation}
where the singular integral satisfies $J(n) = \frac{\Gamma(1+1/2)\Gamma(1+1/4)\Gamma(1+1/6)\Gamma(1+1/8)}{\Gamma(25/24)} \cdot \prod_{k \in \{2,4,6,8\}} (k!)^{-1/k} \cdot n^{1/24}$, and $\mathfrak{S}(n) \ge c_0 > 0$ for all $n \ge 1$.
\end{proposition}

\section{Minor Arc Analysis and the Critical Defect}

\begin{theorem}[Minor Arc Exponent Deficit]
Applying Cauchy--Schwarz and Hua's mean value theorems yields:
\begin{equation}
    \left| \int_{\m} f_2 f_4 f_6 f_8 \, \e(-n\alpha) \, d\alpha \right| \le \|f_8\|_{L^2(\m)} \|f_2 f_4 f_6\|_{L^2([0,1])} \ll N^{23/48 + \varepsilon}.
\end{equation}
Consequently, the uncoupled minor arc bound exceeds the main term $J(n) \asymp N^{1/24} = N^{2/48}$ by an exponent of $\Delta = 21/48 = 7/16 = 0.4375$.
\end{theorem}

\begin{conjecture}[Joint Phase Cancellation on Minor Arcs]
To establish $r(n) > 0$ for all $n \ge n_0$, there must exist $\delta > 0$ such that:
\begin{equation}
    \int_{\m} |f_2(\alpha) f_4(\alpha) f_6(\alpha) f_8(\alpha)| \, d\alpha \ll N^{1/24 - \delta}.
\end{equation}
\end{conjecture}

\end{document}